\section*{الفصل \ref{ch:b3:asymmetric-synthesis} --- التخليق اللاتناظري}

\begin{solution}{exo:b3:asymmetric-synthesis:1}
\omterm{def:b3:asymmetric-synthesis:ee}{زيادة مرآوية} 80\,\%: $\mathrm{er} = 1.8/0.2 = 9$ (90\,:\,10). \omterm{def:b3:asymmetric-synthesis:ee}{ونسبة مرآوية} 99\,:\,1: $\mathrm{ee} = 98/100 = 98$\,\%. \omterm{def:b3:asymmetric-synthesis:ee}{وزيادة مرآوية} 50\,\%: 75\,\% و25\,\%.
\end{solution}

\begin{solution}{exo:b3:asymmetric-synthesis:2}
$(1840 - 60)/(1840 + 60) = 0.937$: \omterm{def:b3:asymmetric-synthesis:ee}{زيادة مرآوية} 93.7\,\%، أي \omterm{def:b3:asymmetric-synthesis:ee}{نسبة مرآوية} $1840/60 = 31$.
\end{solution}

\begin{solution}{exo:b3:asymmetric-synthesis:3}
الأولويات O $>$ الإيثيل $>$ H. وإذا رُسم O إلى الأعلى، والإيثيل في أسفل اليسار وH في أسفل اليمين، كان
الوجه الأمامي \textit{Si}، والوجه الخلفي \textit{Re}. وإضافة مجموعة ميثيل من \omterm{def:b3:asymmetric-synthesis:faces}{الوجه Re}، مع
الأولويات OH $>$ الإيثيل $>$ الميثيل $>$ H في الناتج، تعطي (\textit{R})-بوتان-2-ول.
\end{solution}

\begin{solution}{exo:b3:asymmetric-synthesis:4}
الإيثانول: \omterm{def:b3:asymmetric-synthesis:topicity}{مرآويا الموضع} (يبادل بينهما \omterm{def:b3:point-groups:mirror}{مستوي المرآة} في الجزيء؛ والانزياح
نفسه في مذيب لاكيرالي). بوتان-2-ول: \omterm{def:b3:asymmetric-synthesis:topicity}{لامرآويا الموضع} (بجوار مركز فراغي؛
انزياحان مختلفان).
\end{solution}

\begin{solution}{exo:b3:asymmetric-synthesis:5}
\omterm{def:b3:asymmetric-synthesis:ee}{نسبة مرآوية} 199: $RT\ln 199 = \qty{13.1}{kJ/mol}$ عند \qty{298}{K}، و\qty{8.6}{kJ/mol} عند \qty{195}{K}. \omterm{def:b3:asymmetric-synthesis:ee}{والزيادة المرآوية} 80\,\% عند \qty{25}{\celsius} تقابل \omterm{def:b3:asymmetric-synthesis:ee}{نسبة مرآوية} 9،
$\Delta\Delta G^\ddagger = RT\ln 9 = \qty{5.45}{kJ/mol}$؛ وعند \qty{-78}{\celsius}، $\mathrm{er} = \eu^{5450/(8.314 \times 195.15)} = 28.7$، أي \omterm{def:b3:asymmetric-synthesis:ee}{زيادة مرآوية} 93\,\%.
\end{solution}

\begin{solution}{exo:b3:asymmetric-synthesis:6}
\omterm{def:b3:asymmetric-synthesis:ee}{النقاوة الضوئية} $12.0/15.0 = 80$\,\%: أي 90\,\% من (\textit{S}) و10\,\% من (\textit{R}). وصغر الدوران، والشوائب،
وآثار التركيز والمذيب، وعدم خطية العلاقة بين الدوران
والتركيب تجعلها غير موثوقة؛ أما الكروماتوغرافيا فتفصل
المتماكبين المرآويين وتعدّهما مباشرة.
\end{solution}

\begin{solution}{exo:b3:asymmetric-synthesis:7}
$s = \ln(0.40 \times 0.10)/\ln(0.40 \times 1.90) = -3.22/-0.274 = 12$.
\end{solution}

\begin{solution}{exo:b3:asymmetric-synthesis:8}
مع L-($+$)-طرطرات ثنائي الإيثيل، \babelsublr{(2\textit{S},3\textit{S})}-2،3-إيبوكسي هكسان-1-ول؛ ومع D-($-$)-طرطرات ثنائي الإيثيل،
المتماكب المرآوي، \babelsublr{(2\textit{R},3\textit{R})}.
\end{solution}

\begin{solution}{exo:b3:asymmetric-synthesis:9}
L = الفينيل، M = الميثيل، S = H. ومن أجل الألدهيد (\textit{S}) يكون للمطابق النشط الفينيل
عموديًا على مستوي الكربونيل والهيدروجين بجانب مسار
المحب للنواة، الذي يهاجم في وضع anti من الفينيل. واستخراج التشكيلات يعطي
\babelsublr{(2\textit{S},3\textit{S})}-3-فينيل بوتان-2-ول متماكبًا لامرآويًا رئيسًا.
\end{solution}

\begin{solution}{exo:b3:asymmetric-synthesis:10}
في البداية يوجد المتماكبان المرآويان بكميتين متساويتين، فتتكوّن النواتج
بالنسبة $k_{\mathrm{fast}} : k_{\mathrm{slow}} = s$: $\mathrm{ee}_p = (s - 1)/(s + 1)$، أي 0.905 من أجل $s = 20$. ومع تقدم التفاعل
يُستنفد المتماكب السريع فتنخفض \omterm{def:b3:asymmetric-synthesis:ee}{الزيادة المرآوية} للناتج، بينما ترتفع \omterm{def:b3:asymmetric-synthesis:ee}{الزيادة المرآوية} للركيزة
نحو 1: فيمكن دائمًا تنقية الركيزة بالمضي أبعد؛ أما الناتج
فلا.
\end{solution}

\begin{solution}{exo:b3:asymmetric-synthesis:11}
مع تراسم سريع: مردود حتى 100\,\%، $\mathrm{er} = 99$، أي \omterm{def:b3:asymmetric-synthesis:ee}{زيادة مرآوية} 98\,\%. ودونه، عند تحويل 50\,\%:
مردود 50\,\% على الأكثر، \omterm{def:b3:asymmetric-synthesis:ee}{وزيادة مرآوية} للناتج 93\,\% (\omterm{def:b3:asymmetric-synthesis:ee}{والزيادة المرآوية} للركيزة هي
نفسها، 93\,\%).
\end{solution}

\begin{solution}{exo:b3:asymmetric-synthesis:12}
حين يتحول المعقدان أحدهما إلى الآخر أسرع من تفاعلهما، تكون نسبة النواتج
$k_B[\mathrm B]/(k_A[\mathrm A]) = 600/10 = 60$: فيعطي المعقد الأقل وفرة B نسبة 98\,\% من الناتج (\omterm{def:b3:asymmetric-synthesis:ee}{زيادة مرآوية} 97\,\%). 
والانتقائية يحددها الفرق بين الحالتين الانتقاليتين، لا
أعداد الوسائط.
\end{solution}

\begin{solution}{pb:b3:asymmetric-synthesis:1}
\textbf{1.} لكربون الألكين الحامل للنيتروجين ثلاثة مستبدلات مختلفة
ووجهان مختلفان؛ والهيدروجين المضاف إليه يجعل الكربون $\alpha$ في الحمض
الأميني مركزًا فراغيًا.

\textbf{2.} الخليط الراسيمي: فالوجهان يُهاجمان بالسرعة نفسها.

\textbf{3.} يرتبط أكسجين الكربونيل فيه بالروديوم مع الرابطة C=C: فتُمسك الركيزة
مخلوبةً، في هندسة ثابتة تستطيع الربيطة التمييز فيها.

\textbf{4.} يبني جيبًا كيراليًا حول الفلز (ذا تناظر $C_2$): فطريقتا
ربط الركيزة، بأحد الوجهين أو بالآخر، لامرآويتان، بطاقتين وتفاعليتين
مختلفتين.

\textbf{5.} لا: فالمعقد الأقل وفرة يتفاعل مع الهيدروجين أسرع بكثير (كيرتن--هاميت).

\textbf{6.} يقع الفلز على أحد وجهي الألكين المرتبط، وتُنقل ذرات الهيدروجين
من الفلز.

\textbf{7.} $\mathrm{er} = 1.95/0.05 = 39$.

\textbf{8.} $\Delta\Delta G^\ddagger = RT\ln 39 = 8.314 \times 298.15 \times 3.664 = \qty{9.1}{kJ/mol}$.

\textbf{9.} $\mathrm{er} = \eu^{9082/(8.314 \times 195.15)} = 270$: ee 99.3\,\%.

\textbf{10.} $\mathrm{er} = \eu^{11082/(8.314 \times 298.15)} = 87$: ee 97.7\,\%.

\textbf{11.} أقل من رابطة هيدروجينية نموذجية: فاختلاف صغير في التماسات
الفراغية أو تآثر ضعيف واحد بين الركيزة والربيطة يحسم
النتيجة، ولهذا يصعب تصميم مثل هذه الحفازات انطلاقًا من المبادئ الأولى.

\textbf{12.} تنخفض السرعة، ويتغير ذوبان الهيدروجين وسلوك الحفاز
عند درجة الحرارة المنخفضة، وتبريد مفاعل كبير يكلّف طاقة.

\textbf{13.} 97.5 g من المتماكب المرآوي (\textit{S}) و2.5 g من المتماكب (\textit{R}).

\textbf{14.} $97.5 - 5.0 = \qty{92.5}{g}$ من البلورات، نقية مرآويًا عمليًا.

\textbf{15.} $92.5/100 = 92.5$\,\% من الناتج الخام (95\,\% من المتماكب المرآوي (\textit{S})).

\textbf{16.} تأخذ البلورات المتماكب المرآوي الزائد، ويبقى الجزء الراسيمي في
المحلول مع المتماكب الثانوي؛ أما الخليط الراسيمي فلا زيادة فيه تُجمع، 
وتعطي البلورة بلورات راسيمية.

\textbf{17.} بكروماتوغرافيا HPLC الكيرالية مقارنةً بمرجع راسيمي.

\textbf{18.} 50\,\%: فلا يبقى إلا المتماكب المرآوي البطيء.

\textbf{19.} $s = \ln(0.45 \times 0.05)/\ln(0.45 \times 1.95) = -3.79/-0.131 = 29$.

\textbf{20.} 45\,\% من الخليط الراسيمي (88\,\% من المتماكب المرآوي المطلوب الموجود في البداية،
\omterm{def:b3:asymmetric-synthesis:ee}{بزيادة مرآوية} 95\,\%).

\textbf{21.} مع تراسم سريع للركيزة كان يمكن أن يقترب المردود من 100\,\%.

\textbf{22.} تحوّل كل الطليعة اللاكيرالية الرخيصة إلى المتماكب المرآوي الصحيح،
بحفاز يُستعمل بكميات ضئيلة ودون رمي أي متماكب مرآوي.

\textbf{23.} الروديوم سام وبقاياه في الأدوية محدودة ببضعة أجزاء في
المليون؛ وتُقاس بمطيافية الكتلة بالبلازما المقترنة حثيًا.

\textbf{24.} تقابل \omterm{def:b3:asymmetric-synthesis:ee}{الزيادة المرآوية} 95\,\% عند \qty{25}{\celsius} قيمة $\Delta\Delta G^\ddagger = RT\ln 39 \approx \qty{9.1}{kJ/mol}$.
\end{solution}
